% 中文测试论文——NTex ctex 战役收官验收样张
% 纯中文长文档：多级标题、段落、列表、表格、数学混排
\documentclass[12pt]{article}
\font\zhf=FandolSong-Regular at 12pt
\font\zhhei=FandolSong-Regular at 14.4pt
\font\zhlarge=FandolSong-Regular at 17.28pt
\utfinputmode=1
\begin{document}
\zhf

\title{中文排版测试论文：自然语言处理中的注意力机制研究综述}
\author{测试作者\\测试大学计算机科学系}
\date{2026年9月29日}
\maketitle

\begin{abstract}
本文是一份用于验证中文排版引擎能力的测试文档。全文采用自然语言处理领域的
综述体裁，包含多级章节结构、大量的中文段落文本、数学公式、表格与列表环境，
用于全面检验排版引擎在中文场景下的断行、标点挤压、字体渲染、页码与目录
等核心能力。注意力机制是现代深度学习模型的核心组件之一，自二十世纪九十年代
提出以来不断发展，并在二零一七年之后成为序列建模的主导范式。本文系统梳理了
注意力机制的基本原理、主要变体及其在机器翻译、文本摘要、问答系统等任务中的
应用进展，并讨论了当前面临的挑战与未来可能的研究方向。
\end{abstract}

\section{引言}

注意力机制的思想来源于人类视觉系统：人们在观察复杂场景时，并不会对场景中的
每一个细节给予同等程度的关注，而是会聚焦于少数显著区域。在深度学习领域，
这一思想被形式化为一种可微分的权重分配机制——模型根据当前任务的需要，
动态地计算输入序列中各个位置的重要性权重，并据此对信息进行加权聚合。

早期的注意力机制主要用于对齐源语言与目标语言的单词。传统的编码器—解码器
框架将源句子的全部信息压缩为一个固定长度的向量，这在长句子场景下会造成
明显的信息损失。注意力机制的引入使得解码器在生成每一个目标词时都可以
直接访问源句子的所有位置，从而显著缓解了这一瓶颈问题。

\section{注意力机制的基本原理}

\subsection{查询、键与值}

现代注意力机制的标准形式化将输入分解为三个角色：查询向量、键向量与值向量。
查询代表当前需要处理的位置，键与值代表序列中的所有位置。注意力权重通过
查询与键之间的相似度计算得到，常见的相似度函数包括点积、加性连接与
双线性形式。

缩放点积注意力的计算公式如下：
\begin{equation}
\mathrm{Attention}(Q,K,V)=\mathrm{softmax}\!\left(\frac{QK^{\top}}{\sqrt{d_k}}\right)V
\end{equation}
其中 $Q\in\mathbb{R}^{n\times d_k}$ 表示查询矩阵，$K$ 与 $V$ 分别表示
键矩阵与值矩阵，$d_k$ 为键向量的维度。缩放因子 $\sqrt{d_k}$ 的作用是
防止点积值过大从而导致 softmax 函数进入梯度饱和区域。

\subsection{多头注意力}

多头注意力机制将输入投影到多个不同的子空间，在每个子空间中独立地计算
注意力，最后将各个头的输出拼接起来再做一次线性变换。这种设计使得模型
能够同时关注不同位置、不同层面的信息：有的头专注于捕捉句法结构，
有的头则更关注语义相似性。

\section{主要变体与发展}

\subsection{自注意力与因果掩码}

自注意力机制是指查询、键与值均来自同一序列的注意力形式。在语言模型的
场景下，为了保持自回归性质，需要引入因果掩码——将每个位置对后续位置的
注意力权重置为零。这使得模型在训练时可以并行计算所有位置的损失，
同时保持生成时的合法依赖关系。

\subsection{稀疏注意力与线性注意力}

标准自注意力的计算复杂度与序列长度的平方成正比，这在长文档处理场景下
构成了显著的计算瓶颈。为此，研究者提出了一系列稀疏化与线性化方案：
局部窗口注意力、分块注意力、低秩近似以及核函数方法等。这些方案在
保持模型能力的前提下，将计算复杂度降低到与序列长度线性或近似线性。
各类方法的效果对比见表~\ref{tab:compare}。

\begin{table}[h]
\centering
\caption{主要注意力变体的复杂度与特点对比}
\label{tab:compare}
\begin{tabular}{lll}
\hline
变体 & 复杂度 & 主要特点 \\
\hline
标准自注意力 & $O(n^2 d)$ & 全局交互，表达能力强 \\
局部窗口 & $O(nwd)$ & 计算高效，感受野受限 \\
线性注意力 & $O(nd^2)$ & 核函数近似，长序列友好 \\
稀疏混合 & $O(n\sqrt{n}\,d)$ & 层次化分块，兼顾全局 \\
\hline
\end{tabular}
\end{table}

\section{典型应用}

\subsection{机器翻译}

机器翻译是注意力机制最早取得突破的应用领域。基于注意力的编码器—解码器
模型在多个语言对上大幅超越了传统的统计机器翻译方法，随后出现的
Transformer 架构更是完全抛弃了循环结构，仅依靠堆叠的自注意力层就
达到了当时最先进的翻译质量。

\subsection{文本摘要与问答}

在抽取式摘要任务中，注意力权重可以直接解释为句子的重要性得分；
在生成式摘要中，编码器—解码器注意力则负责源文档与摘要之间的内容选择。
问答系统的匹配模块同样大量采用了注意力机制来建模问题与文档之间的
细粒度交互关系。

\section{挑战与展望}

尽管注意力机制取得了巨大成功，它仍然面临若干挑战：一是超长序列的
计算与显存开销；二是注意力权重可解释性的理论支撑仍然薄弱；三是
位置编码方案的设计缺乏统一的原则。未来的研究方向包括更高效的结构
设计、与检索增强方法的结合，以及在多模态场景下的统一建模。

\section{结论}

本文回顾了注意力机制的基本原理、主要变体与典型应用。从早期的对齐
机制到今天的 Transformer 架构，注意力已经成为深度学习不可或缺的
基础组件。我们期待更高效、更可解释的注意力变体不断涌现，进一步
推动自然语言处理技术的发展。

\end{document}
